\summarychapter{28}{``Another Man'' at Manresa}{The Catholic Reformation}

\begin{inbrief}
Rome was slow to answer the Protestant revolt, but when it did it called up new spiritual warriors,
convened a militant council and reformed the papal office, and by 1600 it had confined Protestantism
to roughly the northern third of Europe. The chapter describes the revival of Catholic piety, the
reforming pope Paul III, Ignatius Loyola and the Jesuits, and the Council of Trent, which gave modern
Roman Catholicism its shape.
\question{How did the Church of Rome respond to the Protestant challenge?}
\end{inbrief}

In 1521, the year of Luther's stand at Worms, a cannonball shattered the leg of a young Spanish
nobleman fighting the French at Pamplona. In his long convalescence he read, out of boredom, a book
of saints' lives and a life of Christ, and his conversion began. At the abbey of Montserrat he
exchanged his gentleman's clothes for a pilgrim's and left his sword and dagger at the shrine of the
Black Virgin. For nearly a year at Manresa, north of Barcelona, he begged, wore a barbed girdle,
fasted and passed through depression to the edge of suicide, until one day by the River Cardoner a
wave of illumination made Ignatius Loyola (1491--1556), in his own words, ``another man.'' Out of that
rebirth he made a plan of spiritual discipline, and out of the plan came the Society of Jesus, the
greatest single force in Rome's campaign to win back what Protestantism had taken. Some call the
movement a Counter Reformation, an attack on Protestantism; others a Catholic Reformation, a genuine
revival. Shelley says it was both: its roots are older than Luther, but its form was largely
determined by the Protestant attack.

\begin{sidebox}{Facets of the Counter Reformation}
(1) Ignatius became a major voice of Catholic spirituality and organized the Jesuit order. (2) Other
spiritual voices emerged, among them Teresa of Avila. (3) Catholics answered Protestant scholarship
with work on the languages, especially in Spain, and with training for priests. (4) The Council of
Trent stated doctrine in reaction to Protestant teaching. (5) The church began a global missionary
expansion. (6) The Inquisition, adapted from Roman legal procedure, proved an effective but brutal
way to detect and stop heresy.
\end{sidebox}

\ssection{Return to the Spiritual}
Mystical experience was a large part of the recovery. The century produced a remarkable variety of
Catholic saints: the English lawyer and statesman Thomas More; Francis of Sales, cheerful missionary
to the Calvinists; Charles Borromeo, the somber reforming archbishop of Milan; the Spanish mystic
Teresa of Avila; and Loyola. Even before Luther's theses a group of aristocrats in Rome had formed
the Oratory of Divine Love, convinced that the reform of church and society begins in the individual
soul. It numbered perhaps fifty but stirred reform in the old orders and supplied leaders: Jacopo
Sadoleto, who debated Calvin; Reginald Pole, who tried to bring England back under Mary; and Gian
Pietro Caraffa, later Pope Paul IV.

Yet through the 1520s and 1530s Rome took no real step. One reason was politics. Luther had called
for a council as early as 1518, and the emperor and the German princes supported the idea, but the
popes remembered Constance and Basel and knew that many Germans wanted a council without the pope.
The popes were also absorbed in worldly affairs. Clement VII (1523--1534) put the Papal States above
everything and allied himself with France against Charles V. The enraged emperor sent troops toward
Rome. Their commanders were killed, and on May 6, 1527, leaderless Spanish and German mercenaries
stormed the city and plundered and murdered for weeks. The pope took refuge in the Castle of
St.~Angelo and then endured half a year of imprisonment. Many saw the sack of Rome as a visitation of
God on a worldly papacy.

Serious reform began with Paul III (1534--1549), an unlikely reformer with three illegitimate sons
and a daughter, whom the sack seems to have sobered. He began with the College of Cardinals,
appointing champions of reform, among them Sadoleto, Pole and Caraffa. Nine of them formed a
commission under Gasparo Contarini, a peacemaker who favored reconciliation with the Protestants and a
return to the faith of the apostles. Its report of 1537, \emph{Advice Concerning the Reform of the
Church}, traced the church's disorder to a papal office grown too secular and demanded an end to
bribery in high places, abuses of indulgences, evasions of church law and prostitution in Rome.

\ssection{Call for a General Council}
Paul acted on some of this, but his chief answer was to call a general council. He and the emperor
agreed on Trent, a town in northern Italy under imperial control. Francis I of France, fearing that a
council would strengthen Charles, did all he could to prevent it and even incited the Turks against
the emperor. Two wars delayed the opening until 1545, almost three decades after Luther's theses. By
then Rome was under a new austerity. The pope instituted the Roman Inquisition and the Index of
prohibited books, which listed all the Reformers' works and Protestant Bibles; in Spain mere
possession of a banned book was long punishable by death. The Index was kept up to date until 1959
and abolished by Paul VI.

\ssection{Ignatius Loyola}
This militant Rome and this pope approved the Society of Jesus, whose members promised to go wherever
the pope sent them, ``to the Turks or to the New World or to the Lutherans.'' As a boy Ignatius had
left the castle of Loyola near the Pyrenees for a noble's court and grown into an engaging playboy
fond of military games and romances of chivalry. Manresa changed that. Luther came out of his struggle
convinced that the will is enslaved and God alone must deliver; Loyola came out of his believing that
God and Satan stand outside man, who has power to choose between them, and that the disciplined use
of the imagination can strengthen the will to choose God. One exercise makes hell real to every
sense: hear the shrieks, smell the sulphur, taste the tears, feel the flames. The same method can
picture the Nativity or the glories of heaven.

For Ignatius surrender to God meant more schooling. He sat with boys half his age in Barcelona to
learn Latin and studied at Alcal\'a, where he concluded that learning must be organized to be useful,
the seed of the Jesuits' famous plan of studies. His zeal brought him more than once before the
Inquisition. He then spent seven years at the University of Paris, became ``Master Ignatius'' and
gathered his first companions: Peter Faber, Diego Laynez, Alfonso Salmeron, Simon Rodriguez, Nicholas
Bobadilla and, above all, Francis Xavier. With them he shared his \emph{Spiritual Exercises}, four
``weeks'' of meditation that move from sin, death, judgment and hell to Christ's life, death and
resurrection. They prescribe rigorous examination of conscience, penance, and then a resolute
forgetting of guilt once God's forgiveness has been met. They became the basis of every Jesuit's
spirituality and were later used for ordination candidates and lay retreats. James Joyce's Stephen
Dedalus, terrified by a sermon on hell, finds after confession that ``the past was past.'' So
Ignatius intended: ``Send no one away dejected. God asks nothing impossible.'' His followers became
apostles of the possible.

\ssection{The Society of Jesus}
Paul III approved the order in 1540. Its members were to be soldiers of Jesus, mobile and versatile.
To vows of poverty and chastity they added obedience to their superiors and a fourth vow of special
loyalty to the pope, and they were commanded by a general elected for life, the first being Ignatius.
Their aim was to restore the Roman church to the spiritual power and worldly influence it had held
under Innocent III, for Ignatius believed that the living Christ dwells exclusively in the
institutional church. Their most striking trait was the risky attempt to live energetically in the
world without being of it, all things to all men. They have had defenders who give them the highest
places in heaven and critics, such as John Adams writing to Jefferson in 1816, who thought them
deserving of perdition. The first generation set out to convert the heathen and reconvert Protestant
Europe. Francis Xavier went from India to Southeast Asia to Japan, which had never heard the
Christian message. More than anyone else the Jesuits stemmed and sometimes reversed Protestantism in
France, the Low Countries and Central Europe. When Ignatius died in 1556 the order numbered nearly a
thousand and worked on four continents.

No mission mattered more than the one to the Council of Trent. It opened with only thirty-one
fathers and three papal legates, yet became the most important council between Nicaea and Vatican II
(1962--1965). The Jesuits Laynez and Salmeron steered its agenda. It met in three periods, 1545--1547,
1551--1552 and 1562--1563. Italians dominated; France was underrepresented; attendance was never
large. Some Protestants attended the second period, with no result.

\ssection{The Shape of Modern Catholicism}
Trent rejected, almost violently, everything the Reformation stood for. Against justification by
faith alone it refused even a compromise that would have combined an external righteousness credited
to the believer with an internal one infused in him, and tied justification firmly to the process of
becoming righteous. Against grace alone it taught grace together with human cooperation; as Ignatius
put it, pray as though everything depended on God, act as though it depended on you. Against
Scripture alone it asserted the teaching office of the Roman church, popes and bishops, as the
necessary interpreter of the Bible. The pope, the seven sacraments, the sacrifice of the mass,
saints, confession and indulgences all remained. The council narrowed the options that medieval
Catholics had had, in an unmistakably anti-Protestant tone.

Europe only slowly grasped that its religious unity was gone for good. Shelley distinguishes three
stages. At first each side believed that it was the one true church and that the other must soon
yield or wither; that was also the spirit of Trent. In the second stage people admitted, bitterly and
without accepting it, that there was a stalemate. Most still held that truth lay wholly on one side
and that error damned souls and destroyed society, so they fought it with inquisitions, civil wars
and persecutions. Calvinists and Jesuits, each a militant organization claiming a loyalty above
nation, led the fight. The third stage, toleration resting on full acceptance of religious diversity
within a nation, appeared before 1600 only as hints among mystics, Christian humanists such as
Erasmus, radical Protestants such as the Anabaptists and practical politicians such as Queen
Elizabeth. Political and economic forces pushed the same way, as small lordships merged into
nation-states that needed Protestant and Catholic citizens to live together. Some already proposed a
general idea of God as the public standard, with distinctive beliefs tolerated as private devotion.
It was a whisper, but it was the voice of the future, and no Christians resisted it more stubbornly
than the Jesuits.

\begin{keyfacts}
\keypeople{Ignatius Loyola; Francis Xavier; Teresa of Avila; Thomas More; Francis of Sales; Charles
Borromeo; Clement VII; Paul III; Gasparo Contarini; Sadoleto; Reginald Pole; Caraffa (Paul IV); Diego
Laynez; Alfonso Salmeron; Charles V; Francis I.}
\keydates{%
  \kd{1521}{Ignatius Loyola wounded at Pamplona; his conversion begins}
  \kd{1527}{Sack of Rome by the emperor's troops (May 6)}
  \kd{1534--1549}{Paul III, the first reforming pope}
  \kd{1537}{Report of the papal reform commission under Contarini}
  \kd{1540}{Paul III approves the Society of Jesus}
  \kd{1545--1563}{Council of Trent, in three periods}
  \kd{1556}{Death of Ignatius Loyola}}
\keyterms{Catholic Reformation and Counter Reformation; Oratory of Divine Love; Society of Jesus
(Jesuits); \emph{Spiritual Exercises}; fourth vow; Roman Inquisition; Index of prohibited books;
Council of Trent; cooperation with grace; teaching office; toleration.}
\end{keyfacts}
